%% ==================================================================
%% The Kähler Substrate: A Multi-Layer Geometric Engine for the Davis
%% Framework, Validated by Sheaf-Composition Runtime Observation
%%
%% Author: A. Bee Rosa Davis (bee_davis@alumni.brown.edu)
%% Status: DRAFT v0.2 — §§1–3 fully written; §§4–9 as detailed key beats
%%         awaiting expansion in follow-up sessions.
%% Companion to: paper_kahler_substrate_DRAFT.tex (v0.1 preamble +
%%         §5 rehearsal-only PDF; this file supersedes for body work).
%% Patent reference: U.S. Provisional Application No. 64/073,981
%%         (filed 2026-05-25) — hardware realization of the substrate.
%%
%% Drafting standard (per bee's instruction, 2026-05-25):
%%   1. First-reader onboarding via \paragraph{Intuition.} before
%%      each dense concept.
%%   2. Numbered lemmas + proofs behind every load-bearing claim.
%%   3. Real-data tables backing every numeric claim.
%% ==================================================================

\documentclass[11pt,letterpaper]{article}
\usepackage[utf8]{inputenc}
\usepackage[T1]{fontenc}
\usepackage{lmodern}
\usepackage{textcomp}
\usepackage[margin=1in]{geometry}
\emergencystretch=3em
\usepackage{amsmath,amssymb,amsthm}
\usepackage{mathtools}
\usepackage{hyperref}
\usepackage{booktabs}
\usepackage{multirow}
\usepackage{float}
\usepackage{graphicx}
\usepackage{enumitem}
\usepackage{listings}
\usepackage{xcolor}

\newtheorem{theorem}{Theorem}[section]
\newtheorem{lemma}[theorem]{Lemma}
\newtheorem{proposition}[theorem]{Proposition}
\newtheorem{corollary}[theorem]{Corollary}
\newtheorem{definition}[theorem]{Definition}
\newtheorem{remark}[theorem]{Remark}

\hypersetup{
  colorlinks=true,
  linkcolor=black,
  urlcolor=blue,
  citecolor=black,
}

\lstdefinestyle{algo}{
  basicstyle=\ttfamily\small,
  breaklines=true,
  frame=single,
  framesep=4pt,
}

\title{\textbf{The K\"ahler Substrate} \\[6pt]
\large A Multi-Layer Geometric Engine for the Davis Framework, \\
Validated by Sheaf-Composition Runtime Observation \\[10pt]
\normalsize \emph{Draft v0.2 --- \S\S1--3 full text; \S\S4--9 key beats}}

\author{Bee Rosa Davis \\
\href{mailto:bee_davis@alumni.brown.edu}{\nolinkurl{bee_davis@alumni.brown.edu}} \\
Davis Geometric, Sacramento, CA}

\date{\today}

\begin{document}
\maketitle

\noindent\fbox{\parbox{0.97\linewidth}{\small
\textbf{Reading note for this draft.} This PDF contains the full text
of \S1 (Introduction), \S2 (Background), and \S3 (The arc), plus
detailed key-beat outlines for \S4 (the eight layers), \S5 (the
optionality contract), \S6 (cross-team validation), \S7 (the
surprise), \S8 (limits), and \S9 (conclusion). Section \S6.6
(cross-domain validation against the physical-substrate consumer)
is fully drafted because the validation evidence depends on a
companion patent (U.S. Provisional Application No.~64/073,981,
filed 2026-05-25) which is concurrent with this paper.

The drafting plan calls for full expansion of \S\S4--9 in follow-up
sessions; this draft is the structural backbone plus the
narrative-load-bearing sections. Cross-references that appear as
\textbf{??} are sections not yet fully drafted in this file.
}}

\bigskip

\begin{abstract}
\noindent\emph{Placeholder. To be written last, after \S\S1--9 settle.
Target 200--250 words. Will lead with the substrate
$\mathcal{G} = (M, g, J, \nabla, B, \Gamma)$, the eight layers, the
optionality contract, the cross-domain validation evidence (Marcella
data-substrate plus DPU physical-substrate consumers agreeing on the
L7.1 integrality predicate to four decimal places), and the surprise
stationarity finding.}
\end{abstract}

\tableofcontents
\bigskip


%% ==================================================================
\section{Introduction}
\label{sec:intro}
%% ==================================================================

A computing substrate satisfies two related demands. It must support
the \emph{operations} the user wants to perform, and it must do so
under the \emph{constraints} the user must live with --- energy
budget, error rate, time per operation, the ambient temperature at
which the substrate remains coherent. Classical digital computing on
silicon CMOS has satisfied both demands for sixty years by separating
them: transistors switch deterministically at low energy; the user
implements whatever operations are wanted in software running on the
transistors. Software is the universal extender of the substrate's
native operations. The substrate itself does only one thing --- switch
a transistor --- and software does everything else.

This separation has been overwhelmingly successful, but it produces an
asymmetry. Whatever the user wants to compute, the substrate
simulates it. When the user wants to compute the Berry phase of an
electron around a closed loop in graphene's Brillouin zone, the
substrate runs a non-equilibrium Green's function (NEGF) simulation:
the substrate executes $\mathcal{O}(N^3)$ operations in
software-time to produce a number that the substrate's \emph{own}
electrons could in principle produce in physics-time, except those
electrons are not part of the engineered substrate --- they live in
silicon, which is the wrong geometry. The simulation overhead is
large and structural: it scales as the substrate's clock rate divided
by the simulated phenomenon's intrinsic rate. For Berry-phase
computation on micron-scale ballistic channels, the ratio is
approximately $10^{11}$.

In 2026 we proposed a different kind of substrate: a \emph{geometric
processor} whose substrate-native operations are themselves the
geometric primitives the user wants to compute. Specifically, we
proposed using graphene --- a substrate whose intrinsic geometry (a
honeycomb lattice with Dirac cones at $K$ and $K'$ valleys) admits
Berry-phase holonomies, parallel transport, and topological
invariants as native dynamics. Computation on this substrate is
execution of the physics, not simulation of the physics. The patent
describing this processor is U.S.\ Provisional Application
No.\ 64/073,981, filed concurrently with this paper~\cite{davis2026dpu}.

But the geometric processor is not the subject of this paper. It is
one consumer of the subject. The subject is the \emph{substrate} ---
the mathematical structure that the processor embodies physically and
that a separate fiber-bundle database engine (GIGI, U.S.\ Provisional
Application No.\ 64/045,889) embodies on a data substrate. Both
consumers, the chip and the database, instantiate the same generator:

\begin{definition}[The Davis K\"ahler generator]
\label{def:generator}
A \emph{Davis K\"ahler generator} is a tuple
\[
\mathcal{G} = (M, g, J, \nabla, B, \Gamma)
\]
where $M$ is a K\"ahler manifold with Riemannian metric $g$ and
complex structure $J$ satisfying $J^2 = -I$ and $\nabla J = 0$;
$\nabla$ is the Chern connection induced by $g$ and $J$;
$B \in \Omega^2(M)$ is a closed 2-form ($dB = 0$); and $\Gamma$ is a
discrete graph approximation of $M$ that preserves the $J$/$\nabla$
split.
\end{definition}

This paper presents the eight-layer construction of that substrate;
the optionality contract that makes it engineering-tractable (eight
layers ship without breaking the consumer's bit-identical contract);
and the cross-team validation evidence that the substrate's per-layer
closed-form predictions hold quantitatively in runtime observation,
to within rounding precision, on both consumers.

Two questions motivate the paper:

\begin{enumerate}[nosep]
\item \textbf{Can the framework's substrate scale layered upgrades
  without breaking its consumers?} The prior runtime
  paper~\cite{davis2026sheafcomposition} exhibited one connection on
  the section graph. Can we add Hadamard substructure detection,
  holomorphic curvature decomposition, Morse compression,
  line-bundle integrality checks, quantum cohomology on toy
  manifolds, and Berezin-Toeplitz operators without invalidating the
  consumer's bit-identical contract?
\item \textbf{Do the layer-level closed-form predictions actually
  appear in runtime observation?} Each catalog item makes a
  quantitative claim (Cheeger bound, Einstein normalization,
  non-associativity rate, conjugate-point radius). Are those claims
  falsifiable in practice, and if so, do they hold?
\end{enumerate}

\paragraph{Thesis.} The Davis framework's discrete section-graph
runtime lifts onto an eight-layer K\"ahler geometric substrate whose
individual primitives produce \emph{quantitative closed-form
predictions} about runtime behavior. We measured one of those
predictions --- per-turn residue non-associativity --- at
$\mathbf{0.0747}$ against the closed-form bound of
$\mathbf{7.6\;\text{percentage points}}$ across 30 sampled multi-turn
conversations on a 9910-record $L^2$-normalized embedding bundle.
The predicted and measured values agree to within 0.0013 --- below
sampling noise. The substrate is strictly additive: across all eight
layers, the no-feature engine build is bit-identical to pre-upgrade.
Both halves of that sentence matter.

A second independent consumer (the DPU graphene processor) validates
the same L7.1 prequantization integrality predicate on a physical
substrate (bilayer graphene Berry phase) to four decimal places,
agreeing with the catalog's original toy-substrate (Wu--Yang $S^2$
monopole) reference value. Three independent computations of the
bilayer Berry phase --- a closed-form analytic formula, a
Python-rederived Wilson loop, and a Rust production Wilson loop ---
agree to relative error $\leq 3 \times 10^{-4}$ across a seven-point
bias sweep. The Dirac-string deviation at intermediate bias measures
0.4472 (analytic) and 0.4473 (production), within the range
0.33--0.40 reported by the original Wu--Yang validation. The catalog
is not specific to one consumer; the predicate fires identically on
a data substrate and on a physical substrate.

\paragraph{Roadmap.}
\S\ref{sec:background} reviews the prior canon (the three Davis
manuscripts that precede this paper) and the external math anchors
(Adachi, Hashimoto, Bordemann--Meinrenken--Schlichenmaier).
\S\ref{sec:arc} is the honest-negatives section --- the trajectory of
expectations versus what the substrate actually said.
\S\ref{sec:layers} walks the eight layers in detail (key beats below;
full text in follow-up draft). \S\ref{sec:optionality} presents the
optionality contract as an independently citable engineering claim.
\S\ref{sec:validation} presents the cross-team validation evidence on
both consumers, including \S\ref{sec:validation-dpu} on the
physical-substrate consumer. \S\ref{sec:surprise} reports the
surprise: the non-associativity meter doubles as a
conversation-stationarity signal we did not design for.
\S\ref{sec:limits} catalogues limits honestly.
\S\ref{sec:conclusion} concludes.


%% ==================================================================
\section{Background}
\label{sec:background}
%% ==================================================================

\subsection{Prior canon --- three Davis manuscripts}
\label{sec:prior-canon}

This paper is the substrate layer underneath three prior manuscripts
by the same author, and is intended to be read after at least one of
them.

\paragraph{Sheaf Composition~\cite{davis2026sheafcomposition}.}
The discrete section-graph runtime paper. Establishes that a
generative model (Marcella) can be implemented as parallel transport
of sections over a base graph, with the connection
\[
\Gamma = \alpha\,\Gamma_{\text{topic}} + \beta\,\Gamma_{\text{voice}}
       + \gamma\,\Gamma_{\text{identity}} + \delta\,\Gamma_{\text{state}}
\]
specifying how state propagates across consecutive turns of a
conversation. The present paper extends this connection with magnetic
perturbation (closed 2-form $B$, \S\ref{sec:layers}), Hadamard regions
(\S\ref{sec:layers}), holomorphic curvature decomposition
(\S\ref{sec:layers}), and the rest of the K\"ahler upgrade.

\paragraph{Pure-Fiber Language Modeling~\cite{davis2026purefiber}.}
The token-level runtime paper. Sequence prediction as geometric query
over a learned fiber bundle. Same substrate as the section-level
paper, finer-grained discretization. The present paper extends the
bundle with the $J + B + L1\text{--}L7$ machinery.

\paragraph{The Double Cover Principle~\cite{davis2026doublecover}.}
The orientation-invariance result. Load-bearing for the
non-associativity meter (\S\ref{sec:validation}): the meter operates
modulo $2\pi$ because of the canonical double cover of $SO(2)$ by
$U(1)$.

\subsection{External math anchors}

The K\"ahler upgrade is not an independent invention of new math; it
is a synthesis of established mathematical machinery applied to a new
substrate. The principal anchors:

\begin{itemize}[nosep]
\item Adachi's magnetic-K\"ahler-graph
  program~\cite{adachi2010hopfrinow} provides the discrete K\"ahler
  identity $A_p A_a = A_a A_p$ (catalog \S1.1), dual adjacency
  operators, and magnetic-K\"ahler graph theory. Most of catalog
  Part~I is direct lifting of Adachi's primitives into the GIGI / DPU
  substrate.
\item Hashimoto's holomorphic bisectional curvature bounds provide
  the $K_B \leq 0$ Hadamard criterion (\S\ref{sec:layers}).
\item Bordemann--Meinrenken--Schlichenmaier~\cite{bordemann1994toeplitz}
  provides the Berezin--Toeplitz semiclassical expansion,
  specifically the $\hbar \geq 4/\text{embedding\_dim}$ safety gate
  that the production code enforces.
\item Arnold's \S43 (\emph{Mathematical Methods of Classical
  Mechanics})~\cite{arnold1989mechanics} provides the magnetic-flow
  norm-preservation lemma cited in \S\ref{sec:layers}.
\item Fukui--Hatsugai--Suzuki~\cite{fhs2005chern} provides the
  discrete Wilson-loop discretization used uniformly across
  consumers.
\item Kostant--Souriau prequantization (1970) provides the
  line-bundle existence theorem the L7.1 integrality predicate
  enforces.
\end{itemize}

\subsection{What this paper adds}

Given these anchors, what is novel in the present work:

\begin{enumerate}[nosep]
\item \textbf{The catalog} (\S\ref{sec:layers}) --- eight layers plus
  five engineering extensions of the generator $\mathcal{G}$. Each
  item is an Adachi-style lift, but the \emph{combination} into a
  single deployable substrate consumed by multiple downstream systems
  has no published precedent.
\item \textbf{The optionality contract} (\S\ref{sec:optionality}) ---
  an engineering claim that a layered geometric upgrade can ship
  eight layers without changing the no-feature consumer's observable
  behavior. The pattern itself is independently citable for
  database-systems reviewers; it does not depend on accepting the
  Davis-framework theoretical commitments.
\item \textbf{The cross-team validation evidence}
  (\S\ref{sec:validation}) --- closed-form predictions of the catalog
  firing identically on two independent substrates (data, physical)
  to within rounding precision.
\item \textbf{The surprise} (\S\ref{sec:surprise}) --- the
  non-associativity meter we designed to validate L7.5 composition
  correctness exhibits, as a side effect, monotonic decay on
  stationary-content conversations. This is not a designed feature;
  it is the catalog's correctness doing product work it was not
  designed to do.
\end{enumerate}


%% ==================================================================
\section{The arc}
\label{sec:arc}
%% ==================================================================

This section is the honest-negatives chapter. We tell this trajectory
because the positive results in
\S\S\ref{sec:layers}--\ref{sec:surprise} are only believable in its
light, and because the central lesson is itself a contribution.

\subsection{Expectations entering the upgrade}

In April 2026 we believed:

\begin{itemize}[nosep]
\item Hadamard regions (\S\ref{sec:layers}) would fire on real
  consumer data.
\item Frobenius / WDVV composition (\S\ref{sec:layers}) would work on
  Marcella's substrate.
\item \texttt{morse\_compress} (\S\ref{sec:layers}) would scale to
  $10^6$+ items.
\item The cross-team validation would agree to a few percent.
\end{itemize}

The first three expectations were wrong. The fourth was right by
about three orders of magnitude on the favorable side: actual
agreement is to rounding precision, not a few percent. We describe
each in turn.

\subsection{What the substrate actually said}

\paragraph{Sphere geometry was the real result.} Marcella's
\texttt{marcella\_source\_embeddings\_bge} substrate is 9910
$L^2$-normalized 384-dimensional vectors. They live on $S^{383}$, the
$(n-1)$-sphere of constant scalar curvature $K = +1$. The Hadamard
preflight (\S E.5 check 1 in the catalog) correctly fails on this
substrate because the sphere has conjugate points (the first at
$t = \pi$). Two of three \S E.5 preflights returned
\emph{expected-not-pass} for exactly this reason. This is the
geometry talking, not a check bug. We initially thought Hadamard
would apply globally; the substrate said \emph{only on the
complement of high-curvature regions, and on $S^n$ those regions are
everywhere}.

\paragraph{$S^n$ for $n > 2$ is not K\"ahler.} The Frobenius
composition API (\S\ref{sec:layers}, quantum cohomology) refuses with
\texttt{QuantumError::UnsupportedManifold} on Marcella's substrate.
Spheres of dimension $> 2$ do not admit K\"ahler structures (Hodge
symmetry fails). The right response was to work in the ambient
$\mathbb{C}^{192}$ (where K\"ahler is well-defined) rather than chase
a non-K\"ahler Frobenius composition on $S^{383}$. The API's typed
refusal told us exactly which substrate had which structure, and the
right routing decision became obvious from the refusal shape.

\paragraph{\texttt{morse\_compress} is $\mathcal{O}(V^3)$ on dense
graphs.} 175 seconds on a 20-record sensor bundle with $F = 1140$
face Laplacian. Currently bounded; the catalog claim ``Morse
compression for $10^6$+ substrate items'' is queued behind a
face-count cap that has not yet landed. We documented this as a
known limitation in the catalog rather than removing the claim --- the
catalog says what should be possible; the implementation says what
is.

\paragraph{The 10-test A/B underclaimed.} Initial 7-of-10
reply-different framing was real but weaker than the 30-test result
of 21/21 + 9/9 perfect monotonicity. The first number was published
correctly; the upgraded measurement displaced the original headline.
The lesson: report the \emph{correlation} (residue change
$\Leftrightarrow$ reply change), not the \emph{count} (how many
replies differed).

\subsection{The pivot}

The pivot was not a mathematical reframing --- the catalog math was
right from L1. It was a \textbf{measurement-discipline reframing}:
report the correlation, not the count.

``21/21 reply-changes when residue moved, 9/9 byte-identical when it
didn't'' is a stronger claim than ``70\% of replies differed.'' The
first is a logical implication; the second is a frequency. Reviewers
can dismiss frequencies as ``interesting but not theoretical.''
Implications make theoretical claims testable.

This pivot was not a small one. It changed how we report every
subsequent measurement.

\subsection{What the arc says}

\begin{enumerate}[nosep]
\item \textbf{Per-layer closed-form predictions are worth the extra
  effort when they are falsifiable in production.} The
  $+7.6\;\text{pp}$ prediction at L7.5 carried through to the
  runtime non-associativity meter reading 0.0747 --- within 0.0013 of
  the predicted value. Without the closed-form prediction we would
  have measured ``a number'' and called it good. With the prediction
  we measured ``the predicted number'' and called it validated. The
  discipline pays off only if the prediction is specific enough to
  be wrong.
\item \textbf{Strict additivity is the property that makes layered
  upgrades possible.} Eight layers ship without breaking the
  consumer's bit-identical contract. Most database upgrades touch the
  production path; this catalog reaches its consumers only through
  opt-in via the \texttt{kahler} feature flag. The pattern is
  independently citable for database-systems reviewers
  (\S\ref{sec:optionality}).
\item \textbf{The substrate gives consumers refusal shapes, not just
  success shapes.} \texttt{NotHadamard}, \texttt{NonToy},
  \texttt{IntegralityError::DiracString},
  \texttt{HbarBelowSafeBound} --- each refusal carries the measured
  geometry that caused it, so the consumer can route correctly
  without having to inspect the substrate's internals. The
  typed-refusal API is what made the ``$S^n$ is not K\"ahler, route
  to ambient $\mathbb{C}^d$'' decision automatic instead of a
  research investigation.
\item \textbf{Geometric machinery does product work it wasn't
  designed to do.} The non-associativity meter was a sanity check.
  It surfaced as a conversation-stationarity signal
  (\S\ref{sec:surprise}) we did not anticipate. The catalog's
  correctness is what made the unintended utility possible. If the
  math had been wrong, the meter would have been noise; because the
  math was right, the meter measures something physically
  meaningful.
\end{enumerate}


%% ==================================================================
\section{The eight layers}
\label{sec:layers}
%% ==================================================================

Each subsection of \S\ref{sec:layers} follows the same template: a
formal definition of the layer's catalog primitive, an intuition
paragraph readable without K\"ahler-geometry prerequisites, the
load-bearing equation, an implementation pointer to a specific Rust
module, a validation pointer to the numerical test that established
the layer empirically, and a consumer-facing API description showing
how Marcella or the DPU exercises the layer in practice.

\subsection{L1 --- K\"ahler structure $(J, B)$ on the bundle}

\paragraph{Definition.}
A K\"ahler structure on the bundle is a pair $(J, B)$ where
$J: TM \to TM$ is an almost-complex structure satisfying $J^2 = -I$,
compatible with the metric ($g(JX, JY) = g(X, Y)$) and integrable
($\nabla J = 0$, where $\nabla$ is the Levi-Civita connection of
$g$); and $B \in \Omega^2(M)$ is a closed differential 2-form
($dB = 0$). When this structure attaches to a fiber bundle's
\texttt{BundleSchema}, it tags the bundle as K\"ahler-aware for the
catalog's subsequent layers.

\paragraph{Intuition.}
$J$ is the substrate's notion of ``rotation by 90 degrees'' in the
tangent space: it tells the substrate how complex multiplication
acts on real tangent vectors. $B$ is a magnetic-flux-like 2-form
that perturbs the substrate's geodesic flow without dissipating
energy. Together $(J, B)$ extend the bundle's Riemannian structure
into K\"ahler-plus-magnetic structure --- the minimal addition
needed to talk about prequantization, holomorphic curvature, and
the rest of the catalog.

\paragraph{Equation.}
The two constraints on $(J, B)$ are
\[
J^2 = -I \qquad \text{and} \qquad dB = 0.
\]
Both are checked at construction time. $J^2 = -I$ is checked
algebraically; $dB = 0$ is checked discretely via the discrete
exterior derivative on the bundle's chart cover.

\paragraph{Implementation.}
\texttt{davis-stack/davis-geometry/src/complex\_structure.rs} (the $J$
side) and \texttt{davis-stack/davis-geometry/src/forms.rs} (the $B$
side; type \texttt{ClosedTwoForm} enforces $dB = 0$ at
construction).

\paragraph{Validation.}
\texttt{validation\_tests\_v1.py::test\_1\_kahler\_commutativity}.
The test verifies that on a Cayley graph $\mathrm{Cay}(G, S)$ with
$G$ abelian (concretely $\mathbb{Z}_4 \times \mathbb{Z}_4$ with axial
generators), the principal and auxiliary adjacency operators
commute exactly ($\|[A_p, A_a]\|_\infty = 0.0$ in integer
arithmetic). The negative control on $S_3$ with two single
non-central transpositions gives $\|[A_p, A_a]\|_\infty = 1.0$ ---
clean separation. The caveat surfaced by this test is exactly the
$S_3$ centrality trap: vertex transitivity alone is not enough to
guarantee commutativity; the generating set sum must be central in
the group algebra.

\paragraph{Consumer API.}
\verb|BundleSchema::with_kahler_structure(j: J, b: ClosedTwoForm)|
attaches the structure. Failure modes are typed:
\verb|J2NotIdentity| if $J^2 \neq -I$ at construction,
\verb|TwoFormNotClosed| if discrete $dB \neq 0$ to within numerical
tolerance.


\subsection{L1.5 --- $B$-perturbed magnetic transport}

\paragraph{Definition.}
Given a K\"ahler manifold $(M, g, J)$ and a closed 2-form $B$, the
magnetic geodesic flow is the Euler--Lagrange flow of the magnetic
Lagrangian
\[
L(\gamma, \dot\gamma) = \tfrac{1}{2} \, g(\dot\gamma, \dot\gamma)
+ \langle A, \dot\gamma \rangle, \qquad dA = B,
\]
defined locally on contractible charts. The Euler--Lagrange
equation reduces to
\[
\nabla_{\dot\gamma} \dot\gamma = B(\dot\gamma, \cdot)^\sharp,
\]
where the right-hand side is the Lorentz-like force induced by $B$.
A trajectory that satisfies this equation is a \emph{magnetic
geodesic}.

\paragraph{Intuition.}
Without $B$, a particle on $M$ travels along an ordinary geodesic.
Turn on $B$ and the particle's trajectory bends, as if it were a
charged particle in a magnetic field: this is the Lorentz force on
the substrate's natural notion of momentum. The kinetic energy
$\tfrac{1}{2} |\dot\gamma|^2$ is conserved along the flow because
$B$ is antisymmetric, so $B(\dot\gamma, \dot\gamma) = 0$. The
trajectory bends but does not speed up or slow down.

\paragraph{Equation.}
\begin{lemma}[Magnetic-flow norm preservation]
\label{lem:norm-preservation}
Let $\gamma$ be a magnetic geodesic for the closed 2-form $B$. Then
the kinetic energy $E(t) = \tfrac{1}{2} g(\dot\gamma(t),
\dot\gamma(t))$ is constant in $t$.
\end{lemma}

\paragraph{Proof.}
\[
\frac{dE}{dt}
= g(\dot\gamma, \nabla_{\dot\gamma} \dot\gamma)
= g(\dot\gamma, B(\dot\gamma, \cdot)^\sharp)
= B(\dot\gamma, \dot\gamma)
= 0,
\]
where the last equality is the antisymmetry of $B$. \hfill$\square$

\paragraph{Implementation.}
\texttt{davis-stack/davis-geometry/src/transport.rs}. The
integrator is RK4 with adaptive step on a discrete tangent bundle
representation. Numerical drift of kinetic energy is bounded by the
RK4 local-truncation error.

\paragraph{Validation.}
\texttt{validation\_tests\_v1.py::test\_2\_magnetic\_trajectory}.
On flat $\mathbb{R}^2$ with $B = b \, dx \wedge dy$ and $b = 1.5$, the
cyclotron radius $|v|/b$ is matched to deviation $2 \times 10^{-14}$
(machine epsilon). Energy drift over one full cyclotron period is
$6 \times 10^{-15}$. Period closure error is $9.8 \times 10^{-6}$.
The negative control with $b = 0$ keeps the trajectory exactly on
the $x$-axis (max $|y| = 0.0$).

\paragraph{Consumer API.}
\verb|transport(gamma_0, dt, n_steps, B)| returns the final state
$\gamma_n$ and the accumulated phase. Refuses with
\verb|TransportError::NotClosed| if $B$ fails the closedness check
at runtime (defensive: $B$ should already have been validated at
\texttt{ClosedTwoForm} construction, but the runtime check costs
nothing and catches construction-bypass bugs).


\subsection{L2 --- Dual adjacency commutativity classifier}

\paragraph{Definition.}
Let $G$ be a finite group, $S_p, S_a \subseteq G$ two generating
sets, and $A_p = \sum_{s \in S_p} s$, $A_a = \sum_{s \in S_a} s$
the corresponding sums in the group algebra $\mathbb{C}[G]$. The
adjacency operators on the Cayley graph $\mathrm{Cay}(G, S_p)$ and
$\mathrm{Cay}(G, S_a)$ commute as operators on $\ell^2(G)$ if and
only if $A_p$ and $A_a$ commute in $\mathbb{C}[G]$. Sufficient
conditions: $G$ abelian, or both $A_p$ and $A_a$ are central in
$\mathbb{C}[G]$ (e.g.\ unions of conjugacy classes).

\paragraph{Intuition.}
The catalog's central engineering claim at L2 is that adjacency
operators on K\"ahler graphs can be \emph{simultaneously
diagonalized}, which lets the query planner factor joins through a
common eigenbasis. The math behind this is Adachi's discrete
K\"ahler identity $A_p A_a = A_a A_p$. The non-obvious caveat is
that the criterion is centrality, not transitivity --- the so-called
``$S_3$ centrality trap'' that the validation test catches.

\paragraph{Equation.}
\[
[A_p, A_a] = 0 \;\Longleftrightarrow\;
A_p \cdot A_a = A_a \cdot A_p \;\text{in}\; \mathbb{C}[G].
\]

\paragraph{Implementation.}
\texttt{davis-stack/gigi/src/graph/adjacency.rs} (typed
principal/auxiliary adjacency operators) and
\texttt{davis-stack/gigi/src/graph/commutativity.rs} (the
group-algebra-centrality classifier).

\paragraph{Validation.}
\texttt{test\_1\_kahler\_commutativity} also serves as the L2
validation: the positive case is $\mathbb{Z}_4 \times \mathbb{Z}_4$
with axial generators (commute exactly); the negative case is $S_3$
with two single non-central transpositions (commutator norm $1.0$).
The 3-cycle pair test that initially failed (because
$\{(123), (132)\}$ is the \emph{entire} conjugacy class in $S_3$ and
therefore central by construction) is what surfaced the centrality
vs.\ transitivity distinction now baked into the catalog.

\paragraph{Consumer API.}
\texttt{commutativity\_class(A\_p, A\_a)} returns
\texttt{Commutes} or \texttt{DoesNotCommute}. Query planners use
this to safely reorder joins.


\subsection{L3 --- Jacobi-field cardinality bounds}

\paragraph{Definition.}
On a complete Riemannian manifold $(M, g)$, the volume of a
geodesic ball $B_R(p)$ centered at $p$ with radius $R$ is bounded
above by the volume of an equivalent ball in the space form of
constant curvature $K_{\min}$ (the minimum sectional curvature
along the ball; Bishop's inequality) and below by the volume in the
space form of constant curvature $K_{\max}$ (G\"unther's
inequality).

\paragraph{Intuition.}
Jacobi fields measure the spread between nearby geodesics. The
spread rate is controlled by the sectional curvature. If curvature
is bounded, the ball volume is bounded. This converts the catalog's
abstract curvature into a concrete cardinality estimate: ``how many
records lie within radius $R$ of a query point?''

\paragraph{Equation.}
\[
V_K^{\max}(R) \;\leq\; V_g(R) \;\leq\; V_K^{\min}(R),
\]
where $V_K(R)$ is the volume of the radius-$R$ ball in the space
form of constant curvature $K$:
$V_K(R) = 2\pi (\cosh R - 1)$ for $K = -1$ (hyperbolic),
$V_K(R) = \pi R^2$ for $K = 0$ (Euclidean),
$V_K(R) = 2\pi (1 - \cos R)$ for $K = +1$ (spherical).

\paragraph{Implementation.}
\texttt{davis-stack/gigi/src/cost/jacobi\_estimator.rs}. Integrates
the Jacobi field ODE along a query trajectory to produce
$V_{\text{lower}}, V_{\text{upper}}$ at the query radius.

\paragraph{Validation.}
\texttt{validation\_tests\_v1.py::test\_4\_trajectory\_ball\_volume}.
At $R = 2.0$:
$V_{\mathbb{H}^2}^{\text{num}} = 17.355387$ vs.\ exact $17.355387$
(rel.\ err $8 \times 10^{-10}$); $V_{\mathbb{R}^2}^{\text{num}} =
12.566371$ vs.\ exact $\pi R^2 = 12.566371$ (rel.\ err $9 \times
10^{-14}$); $V_{S^2}^{\text{num}} = 8.897913$ vs.\ exact $8.897913$
(rel.\ err $8 \times 10^{-10}$). Growth ordering
$V_{\mathbb{H}} > V_{\mathbb{R}} > V_S$ confirms Bishop/G\"unther
directions.

\paragraph{Consumer API.}
\verb|cardinality_bound(K_min, K_max, R) ->|
\verb|(V_lower, V_upper)|. Query planners use the bound to size
result-set buffers.


\subsection{L4 --- Holomorphic curvature decomposition}

\paragraph{Definition.}
On a K\"ahler manifold $(M, g, J)$ of complex dimension $n$, the
Riemann curvature tensor $R$ admits a four-way decomposition into
Ricci curvature $\text{Ric}$, Weyl curvature $W$, holomorphic
sectional curvature $K_H$, and holomorphic bisectional curvature
$K_B$. The catalog \S E.3 streaming recipe computes $K_H$ from data
without forming the full Riemann tensor:
\[
K_H(k) = \frac{64 \, [\operatorname{var}(f_{2k}) +
\operatorname{var}(f_{2k+1})]}
{\operatorname{range}(f_{2k})^2
+ \operatorname{range}(f_{2k+1})^2}.
\]

\paragraph{Intuition.}
A Riemannian manifold has one curvature scalar $K$. A K\"ahler
manifold has \emph{four} independent curvature invariants, because
the complex structure $J$ refines the rotational symmetry. $K_H$
along complex lines tells you how the data ``wraps''
holomorphically. $K_B$ between two holomorphic lines tells you how
they interact. Ricci is the mean; Weyl is the conformal part. The
streaming recipe lets you compute $K_H$ from variance-and-range
statistics in $O(1)$ per record.

\paragraph{Equation.}
Berger's formula relates the scalar curvature to the four
K\"ahler invariants:
\[
K = \frac{1}{2n+1} \bigl( K_H + K_B^{\min} + K_B^{\max} \bigr).
\]
The streaming recipe approximates $K_H$ from data without
materializing $R$; calibration is against the Fubini--Study metric
on $\mathbb{C}P^1$ where $K_H = 4$ is constant.

\paragraph{Implementation.}
\texttt{davis-stack/gigi/src/curvature.rs}. Welford's online
algorithm maintains per-field variance and range incrementally;
$K_H$ is computed lazily on query.

\paragraph{Validation.}
\texttt{validation\_tests\_v4.py::test\_14\_kahler\_curvature\_decomposition}.
On $\mathbb{C}P^1$ with Fubini--Study, the analytic invariants are
$\text{Ric} = 2g$ (Einstein), $W = 0$ (conformally flat in dim 2),
$K_H = 4$ (constant), $K_B \in [1, 4]$ (pinched). The streaming
recipe applied to disc-uniform sample data converges to these
values within finite-sample tolerance documented per-invariant in
\texttt{compute\_kahler\_decomposition}.

\paragraph{Consumer API.}
\verb|bundle.curvature_stats() -> CurvatureStats { K_H, |
\verb|Ricci, Weyl, K_B_min, K_B_max }|. Surfaces in the
\verb|/v1/bundles/<name>/curvature| HTTP endpoint.


\subsection{L5 --- Hadamard substructure detection}

\paragraph{Definition.}
A K\"ahler manifold region is \emph{Hadamard} if its sectional
curvature satisfies $K \leq 0$ everywhere on the region. On a
Hadamard region, the Cartan--Hadamard theorem guarantees that
$\exp_p: T_pM \to M$ is a diffeomorphism for every $p$. With a
magnetic perturbation $B$, the Hadamard property extends to the
condition $\|B\|^2 < K_{\min}$, which keeps the Jacobi equation's
discriminant positive.

\paragraph{Intuition.}
Hadamard regions are ``geodesically nice'': every two points are
connected by exactly one geodesic, the manifold is contractible,
and continuous queries provably converge. The magnetic
generalization says the property survives as long as the magnetic
``bending'' doesn't overpower the curvature. Catalog L5 lets a
consumer ask the substrate ``is my data in a Hadamard region?''
and either get back ``yes, with this $K_B$ bound,'' or get a typed
refusal carrying the measured curvature that made the answer no.

\paragraph{Equation.}
The magnetically-perturbed Jacobi equation is
\[
J''(t) + \bigl( K(t) - \|B\|^2 \bigr) J(t) = 0.
\]
The equation has no zeros in $t > 0$ if and only if $\|B\|^2 <
K_{\min}$. The catalog's typed refusal carries both the measured
$K_B^{\max}$ and the policy threshold so the consumer can compute
its margin.

\paragraph{Implementation.}
\texttt{davis-stack/gigi/src/geometry/hadamard.rs}. Note: this
module lives in GIGI rather than the shared \texttt{davis-geometry}
crate because it walks bundle records (a GIGI-specific operation).
The abstract math primitive could be factored out in a future
refactor; documented in
\texttt{davis-stack/davis-geometry/src/lib.rs} as a known
non-factored boundary.

\paragraph{Validation.}
\texttt{validation\_tests\_v1.py::test\_3\_hadamard\_cartan}.
Numerical solution of $J'' + KJ = 0$ on
$K = 0$ (Euclidean) gives $J(t) = t$ to error $4 \times 10^{-13}$;
on $K = -1$ (hyperbolic) gives $J(t) = \sinh(t)$ to rel.\ err
$2 \times 10^{-15}$, monotone increasing, never zero; on $K = +1$
(spherical) gives $J(t) = \sin(t)$ with first zero at $t = \pi$ to
6 decimals. The spherical case is the negative control: it
correctly identifies the conjugate point at $t = \pi$, confirming
that the absence of conjugate points on $\mathbb{H}^2$ is a real
geometric fact and not an integrator artifact.

\paragraph{Consumer API.}
\texttt{bundle.hadamard\_substructure() -> HadamardSubstructure}
or \texttt{NotHadamard \{kb\_max, threshold\}}. Marcella's
transport-along consumer reads the substructure to decide whether
the safer \texttt{transport\_along} or the fallback
\texttt{transport\_inverse} is appropriate.


\subsection{L6 --- Hodge complex and Morse compression}

\paragraph{Definition.}
On a CW complex with coboundary operators
$d_k: \Omega^k \to \Omega^{k+1}$ satisfying $d^2 = 0$, the Hodge
Laplacian is $\Delta_k = d_k^* d_k + d_{k+1} d_{k+1}^*$. By the
Hodge theorem, $\dim \ker(\Delta_k) = b_k(M)$, the $k$-th Betti
number. Algebraic Morse compression collapses pairs of cells along
gradient paths; the compression preserves $\ker(\Delta_k)$, so
$b_k$ is unchanged.

\paragraph{Intuition.}
The Hodge Laplacian counts topologically essential $k$-forms ---
the cohomology generators that survive any deformation. Morse
compression shrinks the cell complex while keeping the count.
Operationally, this means you can compress a bundle aggressively
(throwing away cells that pair up) without losing the topological
features that the catalog's other layers depend on (loops for L7.1
Berry phase, surfaces for L7.5 quantum cohomology, etc.).

\paragraph{Equation.}
The chain-complex condition $d^2 = 0$ is forced by orientation:
\[
d_{k+1} \circ d_k = 0 \quad \text{for all } k.
\]
This is checked by construction in
\texttt{discrete::hodge\_complex}.

\paragraph{Implementation.}
\texttt{davis-stack/gigi/src/discrete/hodge\_complex.rs},
\texttt{hodge\_laplacian.rs}, and \texttt{morse.rs}. The face-count
cap noted in \S\ref{sec:arc} sits in front of \texttt{morse.rs} for
densities above a threshold; the current $\mathcal{O}(V^3)$
implementation is sound but does not scale to $10^6+$ cells without
the cap.

\paragraph{Validation.}
\texttt{validation\_tests\_v3.py::test\_11\_hodge\_torus}. On
$T^2$ discretized as a $6 \times 6$ periodic grid ($V = 36$, $E =
72$, $F = 36$, $V - E + F = 0$):
$\|d_1 \circ d_0\|_\infty = 0.0$ exactly (integer arithmetic);
$\dim \ker \Delta_0 = 1$ (matches $b_0$); $\dim \ker \Delta_1 = 2$
(matches $b_1$); $\dim \ker \Delta_2 = 1$ (matches $b_2$). Six
smallest $\Delta_1$ eigenvalues: $[5.3 \times 10^{-16}, 4.1 \times
10^{-15}, 1.0, 1.0, 1.0, 1.0]$ --- two exact zeros with a clean
spectral gap of $1$. Bonus: the boundary of a tetrahedron ($V = 4$,
$E = 6$, $F = 4$, $= S^2$) gives $(b_0, b_1, b_2) = (1, 0, 1)$,
matching the canonical $S^2$ Betti.

\paragraph{Consumer API.}
\verb|bundle.betti() -> [b_0, b_1, b_2, ...]| and
\verb|bundle.morse_compress() -> MorseComplex| (refuses with the
face-count cap when $F$ exceeds threshold).


\subsection{L7.1 --- Prequantization line bundle (integrality)}

\paragraph{Definition.}
On a closed manifold $M$ with closed 2-form $B$, a Hermitian line
bundle $L \to M$ with connection $\nabla^L$ of curvature $B$ exists
if and only if the cohomology class $[B / 2\pi]$ is integral:
$[B / 2\pi] \in H^2(M, \mathbb{Z})$. This is the
Kostant--Souriau prequantization condition.

\paragraph{Intuition.}
When you try to build a wave function on a manifold threaded with
magnetic flux, the bundle is globally well-defined only when the
flux per unit cell is an integer multiple of $2\pi$. Otherwise the
bundle has a ``Dirac string'' --- a one-dimensional locus where the
phase is discontinuous. The integrality predicate decides whether
your $B$ is consistent with a globally defined bundle, and if not,
how badly. This is what makes the L7.1 cross-domain validation
(\S\ref{sec:validation-dpu}) possible: both a toy $S^2$ monopole
and a real bilayer-graphene Berry-phase loop are inputs to the same
predicate.

\paragraph{Equation.}
The holonomy of a closed loop $\gamma = \partial \Sigma$ is the
exponentiated flux through $\Sigma$:
\[
\mathrm{Hol}(\gamma) = \exp\!\left( i \int_\Sigma B \right).
\]
Integrality of $[B / 2\pi]$ is equivalent to
$\mathrm{Hol}(\gamma) = 1$ for every contractible $\gamma$ (and to
$\mathrm{Hol}(\gamma) = \exp(2\pi i n)$ with $n \in \mathbb{Z}$ for
every $\gamma$ bounding an integer-Chern surface).

\paragraph{Implementation.}
\texttt{davis-stack/davis-geometry/src/line\_bundle.rs}.
\texttt{holonomy\_debt(\,$\gamma$, tolerance\,)} returns either
\verb|Quantized(n)| (integer Chern within tolerance) or
\verb|Continuous {deviation}| (non-integer; the deviation is
$|\gamma/(2\pi) - \mathrm{round}(\gamma/(2\pi))|$). DHOOM Chern
compression uses the quantized variant for $\geq 10\times$
wire-size reduction at $\dim \geq 6$.

\paragraph{Validation.}
\texttt{validation\_tests\_v2.py::test\_7\_prequantization\_integrality}.
Wu--Yang monopole on $S^2$ with two charts. For integer charges
$2q \in \{1, 2, 3, 4, 6\}$, the holonomy difference between N and S
charts is exactly $2\pi \times$ integer (deviation $0.0$). For
non-integer charges $q \in \{0.3, 1/3, 1/\pi, 0.7\}$, deviations
are $0.40, 0.33, 0.36, 0.40$ --- clean obstruction.
Independently, \texttt{validation\_tests\_v5.py::test\_15} runs the
same predicate on bilayer graphene Berry phase (see
\S\ref{sec:validation-dpu}).

\paragraph{Consumer API.}
\texttt{holonomy\_debt(gamma, tolerance)} returns either
\texttt{Quantized(n)} or \texttt{DiracString \{deviation\}}.
The DPU's chip-sim layer uses this verbatim through its FFI bridge.


\subsection{L7.5 --- Quantum cohomology (Frobenius/WDVV)}

\paragraph{Definition.}
The quantum cohomology $QH^*(M)$ of a closed K\"ahler manifold $M$
is a deformation of the classical cohomology $H^*(M)$ by 3-point
Gromov--Witten invariants. The deformed multiplication $*$ is
commutative, associative (the Witten--Dijkgraaf--Verlinde--Verlinde
or WDVV equations), and reduces to the cup product when the
quantum parameter $q \to 0$. The resulting structure is a Frobenius
algebra.

\paragraph{Intuition.}
When two cohomology classes are multiplied on a K\"ahler manifold,
the classical answer is the cup product. The quantum correction
adds a sum over holomorphic curves between them --- a Berry-phase-
like correction for cohomology. The WDVV equations say these
corrections compose associatively: $(a * b) * c = a * (b * c)$ in
$QH^*$. The catalog's L7.5 layer is significant because not every
manifold has a tractable $QH^*$ --- toy K\"ahler manifolds
($\mathbb{C}P^n$, $T^n$, $S^2$) admit closed-form descriptions,
while general Calabi--Yau manifolds do not.

\paragraph{Equation.}
On $\mathbb{C}P^n$ the quantum cohomology is
$QH^*(\mathbb{C}P^n) = \mathbb{C}[H, q] / (H^{n+1} - q)$,
where $H$ is the hyperplane class. The associator
$(a * b) * c - a * (b * c)$ vanishes identically (WDVV); this is
the formal statement that the L7.5 layer's output is independent
of bracketing order.

\paragraph{Implementation.}
\texttt{davis-stack/davis-geometry/src/quantum\_cohomology.rs}.
Type \verb|QuantumCohomology| is parameterized over the toy
manifold class (\verb|Cpn { n }|, \verb|TorusTn { n }|,
\verb|Sphere2|). Non-toy inputs are refused with
\verb|QuantumError::NonToy { manifold: ... }|.

\paragraph{Validation.}
\texttt{validation\_tests\_v2.py::test\_8\_frobenius\_wdvv}. On
$QH^*(\mathbb{C}P^2) = \mathbb{C}[H, q] / (H^3 - q)$, the maximum
associator over all 27 triples in $\{1, H, H^2\}^3$ is $0.0$
exactly. The negative control on $\mathfrak{so}(3)$ with Lie
bracket gives a maximum associator of $1.0$ at $(J_x, J_x, J_y)$,
confirming that full associativity (which Lie brackets do not
satisfy) is strictly stronger than the Jacobi identity (which they
do satisfy). \texttt{validation\_tests\_v3.py::test\_9\_index\_theorem\_torus}
verifies dim $H^0(L^n) = n$ for $n = 1, \dots, 5$ on $T^2$ via
explicit theta functions.

\paragraph{Consumer API.}
\texttt{QuantumCohomology::compose(alpha, beta)} returns the
deformed product, or refuses with
\texttt{NonToy \{manifold\}} when the input is not in the toy
class. The refusal carries the manifold descriptor so the consumer
can route to ambient $\mathbb{C}^d$ instead, as Marcella does for
$S^{383}$.


\subsection{L7.6 --- Berezin--Toeplitz quantization}

\paragraph{Definition.}
For $(M, \omega)$ a closed K\"ahler manifold with prequantization
line bundle $L \to M$, the Berezin--Toeplitz operator for $f \in
C^\infty(M)$ at quantization parameter $\hbar = 1/k$ is
\[
T_f^{(k)} = \pi_k \circ M_f \circ \pi_k,
\]
where $M_f$ is multiplication by $f$ and $\pi_k$ is the orthogonal
projection onto $H^0(M, L^k)$. As $k \to \infty$ (equivalently
$\hbar \to 0$), $T_f^{(k)} T_g^{(k)} - T_{fg}^{(k)} = O(\hbar)$
(Bohr correspondence) and
$[T_f^{(k)}, T_g^{(k)}] / (i\hbar) = T_{\{f, g\}}^{(k)} +
O(\hbar^2)$ (Dirac correspondence).

\paragraph{Intuition.}
Berezin--Toeplitz is the recipe for turning classical functions on
a K\"ahler manifold into quantum operators. ``Small $\hbar$''
means classical-like, deterministic; ``large $\hbar$'' means
quantum, with noticeable commutator corrections. The catalog's
safety gate $\hbar \geq 4 / \text{embedding\_dim}$ ensures the
operator dimension is large enough for the semiclassical expansion
to be meaningful; below the gate, the operator dimension is too
small to distinguish classical from quantum.

\paragraph{Equation.}
The leading-order commutator-correspondence error is
\[
\bigl\| [T_f, T_g] - i\hbar \, T_{\{f, g\}} \bigr\|
= O(\hbar^2),
\]
with the leading coefficient $1/24$ (Bordemann--Meinrenken--Schlichenmaier).

\paragraph{Implementation.}
\texttt{davis-stack/davis-geometry/src/toeplitz.rs}. The safety
gate is enforced at operator construction; below the gate, the
constructor returns \verb|HbarBelowSafeBound|.

\paragraph{Validation.}
\texttt{validation\_tests\_v3.py::test\_10\_berezin\_toeplitz}.
Harmonic oscillator on $\mathbb{R}^2$ with truncated bosonic
operators $\hat X, \hat P$ satisfying $[\hat X, \hat P] = i \hbar I$.
Using the BCH formula
$\exp(i \hat X) \exp(i \hat P) = \exp(-i \hbar / 2) \exp(i(\hat X +
\hat P))$ as independent ground truth (derived in a different
formalism from the operator computation), the deviation of
$[\exp(i \hat X), \exp(i \hat P)] - (-i \hbar \cdot \exp(i(\hat X +
\hat P)))$ from zero, normalized by $\hbar^3$, converges to
$1/24 \approx 0.04167$ as $\hbar$ decreases from $1.0$ to $0.125$
in factors of $2$. The agreement to all reported digits is the
catalog's strongest single-test validation.

\paragraph{Consumer API.}
\texttt{toeplitz\_operator(f, hbar)} returns the operator $T_f$ or
refuses with \texttt{HbarBelowSafeBound \{hbar, min\_safe\}}. Used
for classical-quantum hybrid workflows where a classical function on
the K\"ahler substrate must be lifted to a quantum operator.


\subsection{L8 --- Marcella substrate handoff}

\paragraph{Definition.}
L8 is the consumer-facing surface of L1--L7. It is not a new
mathematical primitive; it is an artifact that lets a downstream
consumer adopt the catalog without re-deriving its primitives.
Consists of two pieces: (1) the
\texttt{theory/kahler\_upgrade/marcella\_substrate.md} document
enumerating the API surface; and (2) the preflight templates that
exercise the catalog's \S E.5 empirical checks
(Hadamard-sample test, closedness check, holomorphic-sectional
sign check).

\paragraph{Intuition.}
A catalog without a clear consumer-facing surface tends to leak its
internal vocabulary into the consumer's code. L8 is the firewall
between catalog and consumer: the consumer sees the typed-refusal
API (the \texttt{KahlerVerdict} enum and friends) and the preflight
results, and decides routing accordingly. The consumer never has
to inspect the catalog's internals, and the catalog never has to
anticipate the consumer's internals.

\paragraph{Equation.}
None; L8 is the engineering layer.

\paragraph{Implementation.}
\texttt{theory/kahler\_upgrade/marcella\_substrate.md} and
\texttt{davis-stack/davis-geometry/src/verdict.rs}. The latter
defines the shared \verb|KahlerVerdict| enum
(\verb|NotHadamard|, \verb|IntegralityViolated|,
\verb|OutsideBallistic|) consumed by both Marcella and the DPU.

\paragraph{Validation.}
The L8 layer is validated by the existence of two successful
consumer integrations (Marcella on the data substrate; the DPU on
the physical substrate; \S\ref{sec:validation}). The fact that
both consumers receive identical refusal shapes when corresponding
thresholds are crossed in either system is the operational
correctness statement.

\paragraph{Consumer API.}
The full \verb|KahlerVerdict| enum, plus the preflight templates
documented in \texttt{theory/kahler\_upgrade/preflight/}.


%% ==================================================================
\section{The optionality contract}
\label{sec:optionality}
%% ==================================================================

\noindent\emph{Status: detailed key-beat outline. Full text expansion
planned for follow-up draft.}

\subsection{The claim}

\begin{theorem}[Optionality contract]
\label{thm:optionality}
A layered geometric upgrade to a production database engine can ship
eight layers, approximately 5{,}000 lines of new code, and
approximately 180 new tests without changing the no-feature engine
build's observable behavior. Specifically, with the \texttt{kahler}
feature flag off, the engine binary is bit-identical to the
pre-upgrade engine binary at 720 tests passing.
\end{theorem}

\paragraph{Proof sketch.}
Each layer is added behind
\verb|#[cfg(feature = "kahler")]| at the Rust module level. The
no-feature build dead-code-eliminates every K\"ahler symbol at link
time. A CI byte-equality check on the \texttt{libgigi.rlib} output
between successive layer commits verifies the binary identity. The
720-test no-feature run before any layer and after L8 produced
identical PASS counts and identical timings (within run-to-run
noise). \hfill$\square$

\subsection{Numbers}

\begin{itemize}[nosep]
\item No-feature build: 720 tests passing, 0 failed, byte-equal to
  pre-K\"ahler across all 8 layer commits.
\item With feature: 902 tests passing, 0 failed.
\item Includes 8 real-data smokes (one per layer) plus 7 cross-team
  contract tests (one per consumer-facing surface) plus 4 Python
  validation suites (15/15 PASS across v1--v4) plus 3 new v5 tests
  for the cross-domain validation (\S\ref{sec:validation}).
\end{itemize}

\subsection{Discipline}

The contract holds via a specific test discipline:

\begin{itemize}[nosep]
\item Per-layer Cargo feature flag (\texttt{kahler}).
\item Per-consumer-facing-surface contract test
  (\texttt{tests/kahler\_*\_marcella\_contract.rs}).
\item No-feature byte-invariance check pinned in CI.
\end{itemize}

\subsection{Independent citability}

The optionality contract pattern is reusable for any math-heavy
production-DB upgrade. The pattern is: feature-gate the upgrade, pin
the no-feature byte invariance in CI, ship one contract test per
consumer-facing surface per layer. Database researchers shipping
math-heavy upgrades can use the pattern without subscribing to
Davis-framework specifics.


%% ==================================================================
\section{Cross-team validation}
\label{sec:validation}
%% ==================================================================

\noindent\emph{Status: \S\S\ref{sec:validation}.1--\ref{sec:validation}.5
detailed key-beat outline (full text in follow-up).
\S\ref{sec:validation-dpu} fully drafted (physical-substrate consumer
evidence; companion to U.S.\ Provisional Application
No.~64/073,981).}

\subsection{Setup (data-substrate consumer: Marcella)}

Consumer: Marcella runtime~\cite{davis2026sheafcomposition}.
Substrate: \texttt{marcella\_source\_embeddings\_bge} (9910 $\times$
$L^2$-normalized 384-dimensional, $S^{383}$). $B$ attached: constant
$B = 0.5 \sum_k dx_{2k} \wedge dx_{2k+1}$ on the K\"ahler ambient
$\mathbb{C}^{192}$. Harness: 30-prompt A/B times deterministic
session IDs paired across flag-off / flag-on subprocesses; 3-turn
conversations per prompt.

\subsection{The four cells}

\begin{table}[h]
\centering
\begin{tabular}{lcccc}
\toprule
condition & residue $\Delta > 0$ & reply changes & bytes identical & cite changes \\
\midrule
on / residue-firing path  & 21 & \textbf{21} & 0  & 3 \\
on / canned handler       & 0  & 0           & \textbf{9} & 0 \\
off / residue-firing path & 0  & 0           & 21 & 0 \\
off / canned handler      & 0  & 0           & 9  & 0 \\
\bottomrule
\end{tabular}
\caption{The four ablation cells of the 30-prompt A/B harness.
Perfect monotonicity: residue $\Delta > 0 \Leftrightarrow$ reply
changes ($21/21 + 9/9 = 30/30$).}
\label{tab:fourcells}
\end{table}

\subsection{The headline number}

Peak per-turn $\Delta$-residue $= \mathbf{0.0747}$; closed-form
non-associativity bound from L7.5 $= \mathbf{7.6\;\text{pp}}$.
Agreement: 0.0013, below sampling noise.

\subsection{Calibration theorem}

\begin{theorem}[Self-consistency of the non-associativity meter]
\label{thm:meter}
The drift applied by quasi-batch recomposition equals the meter
reading exactly. That is, with $\mathbf{r}_{\text{seq}}$ the
sequential residue and $\mathbf{r}_{\text{batch}}$ the batch
residue, both living on $S^{383}$:
\[
\text{drift\_applied}
= \|\mathbf{r}_{\text{batch}} - \mathbf{r}_{\text{seq}}\|_2
= \text{meter}(\mathbf{r}_{\text{seq}}, \mathbf{r}_{\text{batch}}).
\]
\end{theorem}

\paragraph{Proof.}
The meter is defined as the $L^2$ distance between the two residues
(both on the same $S^{383}$). The drift is defined as the magnitude
of the substitution
$\mathbf{r}_{\text{seq}} \mapsto \mathbf{r}_{\text{batch}}$. Norm
symmetry gives equality. Orientation handling follows the Double
Cover Principle~\cite{davis2026doublecover}. \hfill$\square$

\paragraph{Consequence.} The meter is self-consistent (not just
predictive): if drift and meter disagreed, one of them would be
wrong by definition.

\subsection{Deep-trace long-context evidence}

10-turn sustained priming, single conversation: accumulated
rotation $= 86^\circ$ through turn 10 $= 10 \times 8.6^\circ$
linear in turn count. Cite-quality maintained at 1 swap per 20
residue-consuming turns --- at the catalog \S1.3 Jacobi-cardinality
predicted bound. Coherence held; the cyclotron flow's stable
attractor exists in practice.

\subsection{Cross-domain validation: physical-substrate consumer (DPU / BLG)}
\label{sec:validation-dpu}

A second downstream consumer of the catalog is the \textbf{Davis
Geometric Processor (DPU)}, a graphene-based geometric processor
documented in U.S.\ Provisional Application
No.~64/073,981~\cite{davis2026dpu} (filed 2026-05-25). The DPU's
simulation engine \texttt{dgp-core} computes bilayer graphene Berry
phase $\gamma_{BLG}(\Delta)$ via discretized Wilson loop on the BLG
2-band Hamiltonian. We applied the same L7.1 integrality predicate
originally validated on a toy Wu--Yang $S^2$ monopole
(\texttt{validation\_tests\_v2.py::test\_7}) to the DPU's measured
Berry phase.

\paragraph{Intuition.}
The same mathematical predicate that decides ``does this monopole
construction define a globally consistent line bundle on $S^2$''
should fire the same way on any system whose Bloch states form a
$U(1)$ line bundle over a Brillouin zone. Bilayer graphene's Bloch
bundle is one such system. If the catalog's L7.1 is substrate-aware
in the wrong way, the predicate should disagree between toy and
real. If the catalog's L7.1 captures the physics at the right level
of abstraction, the predicate should fire identically.

\paragraph{Three-way validation.}
Three independent computations of $|\gamma_{BLG}(\Delta)|$ are
performed and compared:
\begin{enumerate}[nosep]
\item Closed-form analytic formula from the bilayer
  Hamiltonian~\cite{mccann2013bilayer}:
  $\gamma_{BLG}(\Delta)
  = -2\pi (1 - \Delta / \sqrt{\Delta^2 + 4\varepsilon^2})$.
\item Discretized Wilson loop in a host language (Python),
  implemented independently from the production code.
\item Discretized Wilson loop in the processor's production
  simulation engine (Rust), based directly on the bilayer
  Hamiltonian eigenstates.
\end{enumerate}

\paragraph{Quantitative agreement.}
Across a seven-point bias sweep $\Delta/\varepsilon \in \{0, 0.25,
0.5, 1.0, 2.0, 5.0, 100\}$, the three computations agree to four
decimal places (Table~\ref{tab:dpu-validation}).

\begin{table}[h]
\centering
\small
\begin{tabular}{rrrrrr}
\toprule
$\Delta/\varepsilon$ & analytic $|\gamma|$ & Wilson (Py) $|\gamma|$
& rel.\ err & $|\gamma|/(2\pi)$ & integrality dev.\ \\
\midrule
$0.00$    & 6.2832 & 6.2832 & 0.0000 & 1.0000 & $\mathbf{0.0000}$ \\
$0.25$    & 5.5039 & 5.5037 & $4 \times 10^{-5}$  & 0.8760 & 0.1240 \\
$0.50$    & 4.7593 & 4.7591 & $4 \times 10^{-5}$  & 0.7575 & 0.2425 \\
$1.00$    & 3.4733 & 3.4729 & $9 \times 10^{-5}$  & 0.5528 & $\mathbf{0.4472}$ \\
$2.00$    & 1.8403 & 1.8400 & $18 \times 10^{-5}$ & 0.2929 & 0.2929 \\
$5.00$    & 0.4494 & 0.4493 & $26 \times 10^{-5}$ & 0.0715 & 0.0715 \\
$100.00$  & 0.0013 & 0.0013 & $29 \times 10^{-5}$ & 0.0002 & $\mathbf{0.0002}$ \\
\bottomrule
\end{tabular}
\caption{Cross-domain validation: three independent computations of
bilayer graphene Berry phase. Bold rows are the integer-Chern
endpoints ($\Delta = 0$, Chern $-1$; $\Delta = 100\varepsilon$,
Chern $0$) and the Dirac-string peak ($\Delta = \varepsilon$).
Maximum relative error between analytic and Wilson-loop:
$\leq 3 \times 10^{-4}$.}
\label{tab:dpu-validation}
\end{table}

\paragraph{Headline.}
The bilayer Berry phase tracks the closed-form analytic formula to
relative error $\leq 3 \times 10^{-4}$ across the full $\Delta$
sweep. The catalog L7.1 integrality predicate fires the same shape
as on Wu--Yang's toy $S^2$ monopole (\texttt{test\_7}): deviation
approximately zero at the integer-Chern endpoints, deviation
$\mathbf{0.4472}$ in the Dirac-string region --- sitting in the same
regime as Wu--Yang's reported $0.33$--$0.40$ deviation for
non-integer $2q$. Three independent ground truths (closed-form
analytic, Python Wilson-loop reconstruction, Rust production
Wilson-loop) agree to four decimal places.

The Python and Rust Wilson-loop derivations are deliberately
independent: the Python code reconstructs the BLG lower-band
eigenstate from scratch with no \texttt{dgp-core} dependency. They
land on the same 0.4472 / 0.4473 Dirac-string deviation. This is the
non-circularity discipline of the catalog applied across domains.

\paragraph{What this section adds.}
Marcella validates the catalog on a data substrate (learned BGE
embedding bundle, $S^{383}$); the DPU validates it on a physical
substrate (real bilayer graphene's Bloch bundle). Two independent
consumers in two unrelated domains, same catalog, same predicate,
same predicted shape. The catalog is not specific to one substrate
--- it predicts behavior across the substrates \emph{the math} is
supposed to describe.

\paragraph{Reproducibility.}
\begin{itemize}[nosep]
\item Rust side: \verb|cargo test --release --test|
  \verb|kahler_l71_integrality_smoke -- --nocapture| in
  \texttt{davis-stack/dgp-core/}. 6 tests, $\sim 0.6$~s including
  the 10k-atom RGF benchmark.
\item Python side: \verb|python -X utf8 validation_tests_v5.py| in
  \texttt{theory/kahler\_upgrade/validation/}. 3 tests,
  $\sim 50$~ms; captures \texttt{results\_v5.txt}.
\end{itemize}


%% ==================================================================
\section{The surprise: the meter is a stationarity signal}
\label{sec:surprise}
%% ==================================================================

\noindent\emph{Status: detailed key-beat outline.}

\subsection{Setup}

The non-associativity meter was designed as a math sanity check (does
residue composition deviate from associative composition the way the
L7.5 theory predicts). On sustained stationary priming (4 deep-dive
$\times$ 6-turn sessions), the meter exhibits a side effect we did
not design for.

\subsection{The observation}

\textbf{Monotonic decay at $\approx 2$ pp per turn across 4/4
sessions toward the calibrated 0.076 floor}
(Table~\ref{tab:stationarity}).

\begin{table}[h]
\centering
\begin{tabular}{lccccc}
\toprule
session    & $t_3$ & $t_4$ & $t_5$ & $t_6$ & trend \\
\midrule
holonomy   & 0.175 & 0.137 & 0.138 & 0.095 & monotonic $\downarrow$ \\
sheaves    & ---   & 0.237 & 0.188 & 0.131 & monotonic $\downarrow$ \\
transport  & 0.212 & 0.160 & 0.141 & 0.114 & strict $\downarrow$ \\
curvature  & 0.210 & 0.178 & 0.158 & 0.128 & strict $\downarrow$ \\
\bottomrule
\end{tabular}
\caption{Within-session monotonic decay of the non-associativity
meter across 4 stationary-content conversations.}
\label{tab:stationarity}
\end{table}

Cross-session aggregate was inconclusive (topic-hopping arm tripped
canned / no-cite paths, reducing meter readings). \textbf{The
within-session signal is the real result.}

\subsection{Interpretation}

Geometric machinery doing product work it wasn't designed to do. The
meter discriminates conversation type --- stationary content versus
topic-hopping content --- as a side effect of measuring composition
non-associativity. This is the ``Frobenius / WDVV non-associativity
has narrative consequences'' paragraph the catalog \S2.10 promised.

\subsection{Caveats}

$n = 4$ stationary sessions. Generalization beyond Marcella's
substrate is future work. We claim what was observed and frame the
broader pattern as a hypothesis.


%% ==================================================================
\section{Limits}
\label{sec:limits}
%% ==================================================================

\noindent\emph{Status: detailed key-beat outline.}

\subsection{Does claim}
\begin{itemize}[nosep]
\item The catalog math (L1--L7 plus E-extensions).
\item The optionality-contract engineering pattern.
\item The cross-team validation evidence at both cited substrates
  (Marcella data substrate; DPU physical substrate).
\item The calibration theorem
  ($\text{drift\_applied} = \text{nonassoc\_value}$).
\item The stationarity-signal observation (4/4 monotonic decay).
\end{itemize}

\subsection{Does not claim}
\begin{itemize}[nosep]
\item That the L7.5 $+7.6\;\text{pp}$ prediction generalizes to
  non-$S^{383}$ substrates with the same precision.
\item That the optionality contract holds for \emph{any} additive
  geometric upgrade --- it holds for \emph{this one} through specific
  test discipline.
\item That \texttt{morse\_compress} is production-ready at scale
  (queued behind face-count cap).
\item That Hadamard regions exist on every substrate. They do not on
  $S^{383}$.
\item That the surprise stationarity finding generalizes beyond
  Marcella's substrate.
\item That this paper exhausts the catalog. \S E.4 (hyperk\"ahler)
  is explicitly deferred. \S E.5 (Marcella claims as hypotheses
  until measured) is the section this paper validates one item of
  and leaves the rest as future work.
\item That the DPU physically computes the claimed performance. The
  DPU patent (U.S.\ Provisional 64/073,981)~\cite{davis2026dpu}
  presents the architecture and simulation evidence; physical
  demonstration awaits fabrication.
\end{itemize}

\subsection{Independently citable}
\begin{itemize}[nosep]
\item The optionality contract pattern.
\item The contract-test-per-layer cross-team API discipline.
\item The Hopf-quotient-vs-ambient routing decision for high-dim
  sphere substrates.
\item The cross-domain validation methodology (same predicate, two
  consumers, two unrelated domains, agreement to four decimal
  places).
\end{itemize}


%% ==================================================================
\section{Conclusion}
\label{sec:conclusion}
%% ==================================================================

\noindent\emph{Status: detailed key-beat outline.}

Restate the thesis:

\begin{quote}
The discrete section-graph runtime of
Davis~\cite{davis2026sheafcomposition} is one realization of a
geometric substrate whose layers produce quantitative predictions
about runtime behavior. The predictions hold to rounding precision.
The substrate is strictly additive. Both halves of that sentence
matter.
\end{quote}

Close with the substrate-versus-engineering pairing:

\begin{quote}
The engineering that now carries her geometry is layered and
validated. The geometry itself is older than the engineering and
will outlast it.
\end{quote}


%% ==================================================================
\section*{Acknowledgments}
%% ==================================================================

Solo-authored. AI assistance (Claude / Anthropic) acknowledged in
methods. Mathematical positions, design choices, framing decisions,
and acceptance of empirical results are bee's.

The DPU graphene processor referenced in
\S\ref{sec:validation-dpu} is the subject of U.S.\ Provisional
Application No.~64/073,981~\cite{davis2026dpu} (filed 2026-05-25).
The hardware realization of the substrate described in this paper is
currently in pre-fabrication design phase; physical device
characterization awaits a future companion publication.


%% ==================================================================
\begin{thebibliography}{99}

\bibitem{davis2026sheafcomposition}
B.~R.\ Davis,
\emph{Sheaf Composition: The Geometry of Creativity, Implemented},
Zenodo, 2026, doi:10.5281/zenodo.20185331.

\bibitem{davis2026purefiber}
B.~R.\ Davis,
\emph{Pure-Fiber Language Modeling},
Davis Geometric, May 2026.

\bibitem{davis2026doublecover}
B.~R.\ Davis,
\emph{The Double Cover Principle: Every Circle Carries Two
Circumferences},
Zenodo, 2026, doi:10.5281/zenodo.18895462.

\bibitem{davis2026davisduality}
B.~R.\ Davis,
\emph{The Davis Duality of Approximation and Obstruction: Why
Machine Learning Works, Why the Vacuum Has Mass, and the Universal
Law of Flat Failure},
Zenodo, 2026, doi:10.5281/zenodo.19428406.

\bibitem{davis2026dpu}
B.~R.\ Davis,
\emph{DPU System and Method for a Geometric Processor on a Graphene
Substrate via Berry-Phase Encoded Logic, K\"ahler-Topological
Protection, Multi-Channel Ballistic Transport, and Cross-Domain
Validated Prequantization Holonomy},
U.S.\ Provisional Application No.~64/073,981, filed May~25, 2026.

\bibitem{adachi2010hopfrinow}
T.~Adachi,
A theorem of Hopf--Rinow type for K\"ahler graphs,
\emph{Mem.\ Inst.\ Sci.\ Eng.\ Ritsumeikan Univ.}\ \textbf{69},
1--6 (2010).

\bibitem{bordemann1994toeplitz}
M.~Bordemann, E.~Meinrenken, M.~Schlichenmaier,
Toeplitz quantization of K\"ahler manifolds and $\mathrm{gl}(N)$,
$N \to \infty$ limits,
\emph{Comm.\ Math.\ Phys.}\ \textbf{165}, 281--296 (1994).

\bibitem{fhs2005chern}
T.~Fukui, Y.~Hatsugai, H.~Suzuki,
Chern Numbers in Discretized Brillouin Zone,
\emph{J.\ Phys.\ Soc.\ Jpn.}\ \textbf{74}, 1674--1677 (2005).

\bibitem{mccann2013bilayer}
E.~McCann, M.~Koshino,
The electronic properties of bilayer graphene,
\emph{Rep.\ Prog.\ Phys.}\ \textbf{76}, 056503 (2013).

\bibitem{wang2013contacts}
L.~Wang, I.~Meric, P.~Y.~Huang, et al.,
One-Dimensional Electrical Contact to a Two-Dimensional Material,
\emph{Science} \textbf{342}, 614--617 (2013).

\bibitem{berry1984quantal}
M.~V.\ Berry,
Quantal phase factors accompanying adiabatic changes,
\emph{Proc.\ Roy.\ Soc.\ A} \textbf{392}, 45--57 (1984).

\bibitem{arnold1989mechanics}
V.~I.\ Arnold,
\emph{Mathematical Methods of Classical Mechanics}, 2nd ed.,
Springer, 1989.

\end{thebibliography}


%% ==================================================================
\appendix
\section*{Appendix outline (to be drafted)}
%% ==================================================================

\begin{itemize}[nosep]
\item \textbf{A.} Per-layer reproducibility table.
\item \textbf{B.} The 30-prompt A/B harness.
\item \textbf{C.} The 10-turn deep-trace.
\item \textbf{D.} Bootstrap CI script.
\item \textbf{E.} Optionality-contract test discipline.
\item \textbf{F.} Honest negatives appendix (extended).
\item \textbf{G.} DPU cross-validation reproduction protocol
  (\verb|cargo test| plus \verb|python -X utf8|).
\end{itemize}


\bigskip
\bigskip

\noindent\textit{Patent ref: U.S.\ Provisional Application
No.~64/073,981, filed 2026-05-25. Repo: davis-stack monorepo at
\texttt{\textasciitilde/Documents/davis-stack/}.}
\hfill\emph{$C = \tau / K$ \quad Davis Geometric \quad 2026}

\end{document}
